\begin{proposition}[The finite period coefficients]\hypertarget{tag-01RV}{}\label{inv:A2.2}
\StatementRecord{inv:A2.2}{proved}{01RV}{\SeriesRef{inv:eq:A-s-branch}, \SeriesRef{inv:eq:A-modular-K}, \SeriesRef{inv:A2.3}}

In the coordinate \eqref{inv:eq:A-s-branch},
\begin{equation}\hypertarget{tag-01T4}{}\label{inv:eq:A-mu-value}
\mu_s(0)=i\,C_4,
\qquad
C_4=\frac{\sqrt{2/3}}{\pi E_4(i)^{1/4}}
=\frac{4\sqrt\pi}{3^{3/4}\Gamma(1/4)^2}>0.
\end{equation}
For $r^3=t$, with $r>0$ on the arc $0<t<1$ and $\tau$ travelling
from $\rho$ to $i$ along the unit circle, set
\[
\lambda=\frac{3e^{5\pi i/6}}{2\pi E_6(\rho)^{1/3}}.
\]
Then $\tau=\rho+\lambda r+O(r^2)$ and
\[
K_r=-\frac{3i\lambda}{2\pi}r^2+O(r^3).
\]
The exact expressions have the rational enclosures
\[
\begin{aligned}
0.2366093140807852&<C_4<0.2366093140807853,\\
0.0801022110875382&<\operatorname{Re}K_2<0.0801022110875383,\\
0.1387410994022233&<\operatorname{Im}K_2<0.1387410994022234.
\end{aligned}
\]
\end{proposition}


\begin{proof}
\begin{enumerate}[label=\arabic*.,leftmargin=*]
\item Ramanujan's differential relations give
$E_6'(i)=-\pi i E_4(i)^2$ and
$E_4'(\rho)=-(2\pi i/3)E_6(\rho)$.
Using $j=1728E_4^3/(E_4^3-E_6^2)$, we obtain
\[
j(\tau)-1728=-1728\pi^2E_4(i)(\tau-i)^2+O((\tau-i)^3).
\]
Since $j=1728(1+s^4)$, the chosen branch is precisely
\eqref{inv:eq:A-s-branch}.  Formula \eqref{inv:eq:A-modular-K} now gives
\[
K_s(0)=-\frac4{\pi^2\sqrt{E_4(i)}}=6\mu_s(0)^2,
\]
because $\tau_s(0)=0$.  This proves nonvanishing and the first
formula in \eqref{inv:eq:A-mu-value}; the sign is fixed by the stated
coordinate convention.  Similarly
$j(\tau)=-512i\pi^3E_6(\rho)(\tau-\rho)^3+\cdots$.
This gives the stated $\lambda$, and substitution into
\eqref{inv:eq:A-modular-K} gives the coefficient $-3i\lambda/(2\pi)$.

\item The real period also has the integral expression
\[
\Omega_R=\frac{\sqrt2}{3^{1/4}}
\int_0^{\pi/2}(1-\tfrac12\sin^2\theta)^{-1/2}\,d\theta,
\qquad C_4=\frac{\sqrt2}{3\Omega_R}.
\]
It follows by the elliptic-integral substitution for the three roots
$\sqrt3,0,-\sqrt3$ of $Y^2=X^3-3X$ and the differential
$dX/(2Y)$. The beta-integral evaluation gives the gamma expression
in \eqref{inv:eq:A-mu-value}.

For the rational enclosures, evaluate the divisor-sum series of
$E_4(i)$ and $E_6(\rho)$ by interval arithmetic with $N=100$ terms.
The elementary inequalities $\sigma_3(n)\le n^4$,
$\sigma_5(n)\le n^6$, $e^{-2\pi}<1/500$ and
$e^{-\pi\sqrt3}<1/200$ give respective tail bounds
\[
250(N+1)^4 500^{-N-1},\qquad
1024(N+1)^6 200^{-N-1}\quad(N=100).
\]
The ratios of successive majorant terms are less than $1/480$
and $1/180$, respectively, which proves these tail estimates.
Outward rounding in the finite sums, together with these tail bounds,
gives the three stated enclosures. The preceding modular calculation
proves the exact nonvanishing and the local orders.
\end{enumerate}
\end{proof}